% !TeX spellcheck = en_US
\addsection{Competitive Play}{\art/haste.png}

\iftoggle{printable}{}{\bigbreak}

\begin{multicols*}{2}
We've compiled some custom rules and best practices for fair competitive play.
They are intended for Clash Scenarios, where they reduce randomness and the luck factor.
Keep in mind that these are just the \textbf{opinions of the authors} of this book.
Enjoy!

\subsection*{Magic Arrows}

Except for the Magic Arrow, each Spell appears exactly twice in the Spell Deck.
After players prepare their starting M\&M Decks, it is recommended to remove Magic Arrows from the Spell Deck until no more than two remain.

\subsection*{\textit{(Optional)} Guaranteed Settlement}

When distributing Far Tiles (II--III) between players, each player's pool of Map Tiles must include one Settlement.

\subsection*{\textit{(Optional)} War Machines as Artifacts}

Before the match, shuffle one Card of each War Machine into the Artifact Deck (or the Major Artifact Deck, if you split it as in Tournament Mode).

\subsection*{Round One Mulligan}

At the start of your turn in the \nth{1} Round, you may once reshuffle your hand back into your M\&M Deck and draw a new starting hand.

\subsection*{Trading Post}

Although rules as written state that you can remove only one Card from your hand while visiting a Trading Post, it is recommended to allow removing as many as you like.

\subsection*{Tier V--VII Combat on Hard$+$\\Difficulty}

Treat Level V--VII Combats as one difficulty step lower when playing on Hard or Impossible.
For example, a Level VII Combat on Hard will be against 2 × \svg{azure} Neutral Units.

\subsection*{Victory Points}

When playing with the optional Victory Points from the Tournament book, the following rules apply (unless overruled by the Scenario):
\begin{itemize}
  \item Faction Towns grant 1 VP.
  \item Possessing the Grail at the end of the match grants 3 VPs. If the Grail was transported to that player's Faction Town, it grants 5 VPs instead.
  \item Capturing a Dragon Utopia grants 3 VPs.
\end{itemize}

\subsection*{Combat with Neutral Units}

During Combat with Neutral Units in Clash Scenarios, if a player decides to deploy their Units only on one half (left or right) of the Combat Board, the player positioning the Neutral Units must do the same.

Placing Neutral Units far apart only makes the other player waste \svg{movement}, which slows the game down.
The exception to this rule is when there are three Neutral \svg{unit_ground} or \svg{unit_flying} Units to position, as all of them must go to the front line.
See examples that follow.

\end{multicols*}

\begin{figure*}[!h]
  \mbox{}
  \hfill
  \begin{minipage}[t]{0.36\textwidth}
    \centering
    \includegraphics[width=\linewidth]{\examples/battle_good.png}
    \caption[good protected]{\textit{Neutral Units are positioned correctly.}}
  \end{minipage}
  \hfill
  \begin{minipage}[t]{0.36\textwidth}
    \centering
    \includegraphics[width=\linewidth]{\examples/battle_bad.png}
    \caption[bad protected]{\textit{Not allowed: the Necropolis Units are on the left half, so the Peasants must start there too.}}
  \end{minipage}
  \hfill
  \mbox{}

  \vspace{0.5em}

  \mbox{}
  \hfill
  \begin{minipage}[t]{0.36\textwidth}
    \centering
    \includegraphics[width=\linewidth]{\examples/battle_exception.png}
    \caption[exception protected]{\textit{\textbf{Exception:} Three non-\svg{unit_ranged} Units must be positioned on the front line.}}
  \end{minipage}
  \hfill
  \begin{minipage}[t]{0.36\textwidth}
    \centering
    \includegraphics[width=\linewidth]{\examples/battle_ranged_bad.png}
    \caption[ranged protected]{\textit{The Boars must occupy one of the green Fields.}}
  \end{minipage}
  \hfill
  \mbox{}
\end{figure*}

